% Drawing 141/01-13 -- folder 141, batch 01, archive page 14 (PDF page 15).
% Transcription: à gauche, une sphère : contour rond, équateur en ellipse ; sur le contour, $i$ en haut, $-i$ en bas, un gros point à gauche et un point à droite, chacun avec un chiffre peu lisible (lus $3$ et $1$) ; un point marqué d'un chiffre (lu $2$) sur l'arc arrière de l'équateur ; près du centre, $0$ entouré d'une petite boucle ; un lacet descend de $i$ vers cette boucle, et une courbe va de $i$ à $-i$ en coupant l'arc avant de l'équateur en un point marqué.
% Kind: sphere with its equator; marked points i, -i, 0, 1, 2, 3 and loops. Method 2 (schematic redraw from the facsimile). Status: partial.
% Notes: Redrawn 2026-11-13 as a sphere with its equator, as the editor asked (#148: « ça devrait être dessiné comme une sphère avec l'équateur »): the round outline, and the equator as one true ellipse between the left and right points on the outline (his two freehand arcs), its front half solid and its back half dashed; the marked point of the back arc and the point of the front arc where the curve from i to -i crosses are put on the ellipse (the back point lies about two units higher on the page, where his back arc bulges more). Kept from the page: i (top) and -i (bottom) on the outline, the heavy point at the left and the point at the right, 0 near the centre inside a small loop, the stroke from i down to that loop, the curve from i to -i through the front point. Digits as before, read 2 (back arc), 1 (right), 3 (left): the editor's answer (#148) asks for the sphere and does not settle them, and the transcription's note reads 1 on the back arc with nothing legible at the left and right; still to confirm. Left out: the faint pencil ovals round 0 and a short tick on the curve. Overlaid on the 400 dpi colour crop: outline, points and curves within a unit or two. Not checked by a reader.
% Coordinates: 1 unit = 1 mm = 1/100 of the crop width (bbox in page fractions: [0.2, 0.165, 0.4, 0.295]).
% Redrawn by Claude (claude-opus-5-5), issues #73 and #148, 2026-11-13; not checked by a human; not compiled (no TeX in the session).
% Rights: reproduces part of the work -- behind RENDER_TRANSCRIPTS, never published by itself.
\documentclass[tikz,border=3pt]{standalone}
\usepackage{amsmath,amssymb}
\usetikzlibrary{patterns,arrows.meta}
\tikzset{pencil/.style={black!45}, struck/.style={black!35}, mark/.style={circle,draw,inner sep=1.3pt}, pt/.style={circle,fill,inner sep=0.9pt}, opt/.style={circle,draw,fill=white,inner sep=1.1pt}}
\begin{document}
\begin{tikzpicture}[x=1mm,y=1mm,line cap=round,line join=round,font=\scriptsize]
\draw (50.5,-47.5) circle[radius=29.6];
\draw plot[smooth, tension=0.55] coordinates {(80,-45.5) (79.7,-47.2) (78.4,-48.9) (76.4,-50.5) (73.5,-52.1) (70,-53.5) (65.8,-54.7) (61.2,-55.7) (56.3,-56.5) (51.1,-57) (46,-57.2) (41,-57.1) (36.3,-56.7) (32,-56) (28.3,-55.1) (25.3,-54) (23,-52.6) (21.6,-51.1) (21,-49.5)};
\draw[dashed] plot[smooth, tension=0.55] coordinates {(21,-49.5) (21.3,-47.8) (22.6,-46.1) (24.6,-44.5) (27.5,-42.9) (31.1,-41.5) (35.2,-40.3) (39.8,-39.3) (44.7,-38.5) (49.9,-38) (55,-37.8) (60,-37.9) (64.7,-38.3) (69,-39) (72.7,-39.9) (75.7,-41) (78,-42.4) (79.4,-43.9) (80,-45.5)};
\node[pt] at (50.5,-18) {};
\node[pt] at (51,-76.8) {};
\node[pt] at (21,-49.5) {};
\node[pt] at (80,-45.5) {};
\node[pt] at (56,-37.8) {};
\node[pt] at (57.3,-56.4) {};
\draw plot[smooth, tension=0.55] coordinates {(50.5,-18) (47.5,-30) (46,-39) (50,-46) (57.3,-56.4) (56.5,-62) (54,-69) (51,-76.8)};
\draw plot[smooth, tension=0.55] coordinates {(50.5,-18) (46.5,-34) (43.5,-45) (42.5,-51)};
\draw (42,-56) ellipse[x radius=4, y radius=5];
\node[inner sep=1pt] at (54,-10) {\small $i$};
\node[inner sep=1pt] at (57,-82) {\small $-i$};
\node[inner sep=1pt] at (62,-34.8) {\small $2$};
\node[inner sep=1pt] at (86,-48) {\small $1$};
\node[inner sep=1pt] at (12,-50.5) {\small $3$};
\node[inner sep=1pt] at (42,-56) {\small $0$};
\end{tikzpicture}
\end{document}
